% generated by roam.py handbook export -- не редактировать, править заметки
% 20261010100652-теорема_о_промежуточном_значении.org (body f883a711e4c1)

\begin{theorem}[{Теорема о промежуточном значении}]\label[theorem]{note:6ADBECA2-157A-4C66-B4C4-F5F2AF0899AE}
  Пусть $f: [a;b] \to \R$ непрерывна на $[a;b]$. Тогда

  \begin{enumerate}
    \item если $f(a) \cdot f(b) < 0$, то $\exists c \in (a;b): f(c) = 0$;
    \item $f$ принимает все значения между $f(a)$ и $f(b)$.
  \end{enumerate}
\end{theorem}

\begin{remark}
  Также называется теоремой Больцано-Коши. Непрерывность нужна во всех точках: $f(x) = -1$ при $x \in [-1;0)$, $f(x) = 1$ при $x \in [0;1]$ меняет знак, но нуля нет.
\end{remark}

\begin{proof}
  \proofpart{Пункт 1}
  НУО $f(a) < 0$, $f(b) > 0$. Положим $a_0 = a$, $b_0 = b$ и будем делить отрезок пополам: $c = \frac{a_n + b_n}{2}$.

  \begin{itemize}
    \item Если $f(c) = 0$~--- доказано.
    \item Если $f(c) > 0$, берём $[a_{n+1}; b_{n+1}] = [a_n; c]$; если $f(c) < 0$~--- $[c; b_n]$.
  \end{itemize}

  На каждом шаге $f(a_n) < 0 < f(b_n)$ и $b_n - a_n = \frac{b - a}{2^n} \to 0$. Отрезки вложены и стягиваются, поэтому $\exists! c$: $a_n \le c \le b_n$ при всех $n$, и $a_n \to c$, $b_n \to c$.
  По непрерывности $f(a_n) \to f(c)$ и $f(a_n) < 0$, откуда $f(c) \le 0$; аналогично $f(b_n) \to f(c)$ даёт $f(c) \ge 0$. Значит $f(c) = 0$, и $c \in (a;b)$, так как $f(a), f(b) \ne 0$.

  \proofpart{Пункт 2}
  Пусть $y$ лежит между $f(a)$ и $f(b)$, НУО $f(a) < y < f(b)$. Функция $g(x) = f(x) - y$ непрерывна, $g(a) < 0 < g(b)$, и по пункту 1 $\exists c \in (a;b): g(c) = 0$, то есть $f(c) = y$.
\end{proof}
